\documentclass[10pt,a4paper]{article}
\usepackage[utf8]{inputenc}
\usepackage[T1]{fontenc}
\usepackage[ngerman]{babel}
\usepackage[margin=1.7cm,top=1.5cm,bottom=1.7cm]{geometry}
\usepackage{titlesec}
\titleformat*{\section}{\normalfont\large\bfseries}
\titleformat*{\subsection}{\normalfont\normalsize\bfseries}
\titlespacing*{\section}{0pt}{8pt}{3pt}
\titlespacing*{\subsection}{0pt}{6pt}{2pt}
\usepackage{amsmath,amssymb}
\usepackage{listings}
\usepackage{xcolor}
\usepackage{tikz}
\usepackage{fancyhdr}
\usepackage{parskip}
\usepackage{needspace}

\newcommand{\mcbox}{\raisebox{-1pt}{\fbox{\rule{0pt}{7pt}\rule{7pt}{0pt}}}\hspace{6pt}}
\newcommand{\antwortlinien}[1]{%
  \par\vspace{2pt}%
  \newcount\zeilen \zeilen=#1
  \loop
    \noindent\rule{\linewidth}{0.4pt}\par\vspace{11pt}%
    \advance\zeilen by -1
  \ifnum\zeilen>0 \repeat
}
\lstset{language=Python,basicstyle=\ttfamily\footnotesize,keywordstyle=\bfseries,
  commentstyle=\itshape\color{gray},showstringspaces=false,frame=single,
  framerule=0.3pt,xleftmargin=8pt,framexleftmargin=6pt,aboveskip=6pt,belowskip=6pt}

\pagestyle{fancy}
\fancyhf{}
\fancyhead[L]{\small MDT5/2 -- Session 4: Probabilistische Verfahren}
\fancyhead[R]{\small Übungsaufgaben zur Klausurvorbereitung}
\fancyfoot[C]{\small Seite \thepage}
\renewcommand{\headrulewidth}{0.4pt}

\begin{document}

\begin{center}
  {\Large\bfseries Übungsaufgaben Session 4}\\[4pt]
  {\large Kapitel 6: 3$\sigma$-Regel, Mahalanobis, Histogramm, KDE, GMM -- Praktikum 04}\\[4pt]
  \small Stil der Übungssammlung MDT5/2 -- generierte Aufgaben, keine Originale
\end{center}

\vspace{4pt}

\section*{Aufgabe 1 -- Multiple Choice mit Begründung}
\emph{Markieren Sie jeweils die einzige richtige von drei Aussagen und begründen Sie Ihre Entscheidung.}

\subsection*{1.1 \ Für Wahrscheinlichkeitsdichtefunktionen (PDFs) gilt:}
\mcbox Eine PDF kann an einzelnen Stellen Werte größer als 1 annehmen, solange die
Gesamtfläche 1 ergibt.

\mcbox PDF-Werte sind direkt Wahrscheinlichkeiten und liegen deshalb immer zwischen 0 und 1.

\mcbox Die Wahrscheinlichkeit für einen exakten Einzelwert einer stetigen Zufallsvariablen
ist gleich dem Funktionswert der PDF.

Begründung:
\antwortlinien{2}

\subsection*{1.2 \ Für Kernel Density Estimation (KDE) gilt:}
\mcbox Die Bandbreite $h$ ist der wichtigste Metaparameter und kann unüberwacht z.\,B.
per GridSearchCV über die mittlere log-Dichte der Trainingspunkte eingestellt werden.

\mcbox KDE kann nur unimodale Verteilungen schätzen.

\mcbox Die Bandbreite spielt kaum eine Rolle, solange die Daten normalverteilt sind.

Begründung:
\antwortlinien{2}

\subsection*{1.3 \ Für den Elliptic Envelope / Mahalanobis-Ansatz gilt:}
\mcbox Er funktioniert nur bei unimodalen (Normal-)Verteilungen der Normaldaten zuverlässig.

\mcbox Er berücksichtigt keine Abhängigkeiten zwischen den Merkmalen.

\mcbox Der Parameter \texttt{contamination} gibt die Anzahl der Cluster an.

Begründung:
\antwortlinien{2}

\subsection*{1.4 \ Für Gaussian Mixture Models (GMMs) gilt:}
\mcbox Das EM-Training ist schnell und konvergiert garantiert im globalen Minimum.

\mcbox Ein GMM schätzt die Dichte als Mischung aus so vielen unimodalen Normalverteilungen,
wie Cluster gewählt wurden; neue Punkte erhalten die berechnete Dichte als Score.

\mcbox Die Anzahl der Komponenten wird beim EM-Algorithmus automatisch mitgelernt.

Begründung:
\antwortlinien{2}

\section*{Aufgabe 2 -- 3$\boldsymbol{\sigma}$-Daumenregel rechnen}

Für ein skalares Merkmal wurden aus den normalen Trainingsdaten
$\hat\mu = 10$ und $\hat\sigma = 2$ geschätzt.

a) Beurteilen Sie mit der Daumenregel aus der Vorlesung die neuen Messwerte
$x_1 = 13$, $x_2 = 17$ und $x_3 = 3{,}5$.
\antwortlinien{3}

b) Die Daumenregel wird nun auf jedes Merkmal eines 2D-Problems einzeln angewendet.
Welche implizite Annahme über die Merkmale steckt dahinter, und welche geometrische Form
hat der resultierende Normalbereich?
\antwortlinien{3}

c) Skizzieren Sie 2D-Normaldaten mit stark korrelierten Merkmalen und zeichnen Sie ein,
warum die achsenparallele Box aus b) dann Anomalien \glqq durchlässt\grqq{} --
und welche Form der Mahalanobis-Ansatz stattdessen liefert.

\vspace{5cm}

\section*{Aufgabe 3 -- Mahalanobis-Abstand}

a) Geben Sie die Formel des Mahalanobis-Abstands an und erklären Sie die intuitive
Interpretation (\glqq Abstand in \dots\grqq).
\antwortlinien{3}

b) Warum entsteht der Mahalanobis-Abstand auf natürliche Weise beim Logarithmieren der
multivariaten Normalverteilungs-PDF?
\antwortlinien{3}

c) scikit-learn bietet \texttt{EmpiricalCovariance} und \texttt{MinCovDet} zur Schätzung
der Kovarianzmatrix. Worin unterscheiden sie sich?
\antwortlinien{2}

d) Wozu dient beim \texttt{EllipticEnvelope} der Parameter \texttt{contamination}
und warum verbessert er die Anpassung gegenüber der starren 3$\sigma$-Grenze?
\antwortlinien{3}

\section*{Aufgabe 4 -- KDE verstehen}

Die KDE-Schätzung lautet
$KDE(x) = \frac{1}{Nh}\sum_{i=1}^{N} k\!\left(\frac{x - t_i}{h}\right)$
mit Gauß-Kernel $k$.

a) Erklären Sie die Rolle der Trainingspunkte $t_i$, der Bandbreite $h$ und der
Normalisierung $\frac{1}{Nh}$.
\antwortlinien{3}

b) Was passiert qualitativ mit der geschätzten Dichte, wenn $h$ viel zu klein bzw.
viel zu groß gewählt wird? (Stichwort Overfitting/Underfitting)
\antwortlinien{3}

c) Wie werden bei \texttt{sklearn.neighbors.KernelDensity} Anomalie-Scores für neue
Punkte gewonnen (welche Methode, welche Größe)?
\antwortlinien{2}

\section*{Aufgabe 5 -- Verfahrenswahl zu Plots}

Skizziert sind normale Trainingsdaten zweier Szenarien:

\begin{center}
\begin{tikzpicture}[scale=0.62]
% Szenario 1: banane
\draw[->] (-0.5,0) -- (7,0) node[right] {\small $x_1$};
\draw[->] (0,-0.5) -- (0,5) node[above] {\small $x_2$};
\node at (3.5,-1.2) {\small Szenario 1};
\foreach \a in {10,25,...,170} {
  \fill ({3+2.6*cos(\a)},{0.6+2.9*sin(\a)}) circle (2pt);
  \fill ({3+2.2*cos(\a)},{0.6+2.4*sin(\a)}) circle (2pt);
}
\begin{scope}[xshift=9.5cm]
% Szenario 2: zwei kompakte Cluster
\draw[->] (-0.5,0) -- (7,0) node[right] {\small $x_1$};
\draw[->] (0,-0.5) -- (0,5) node[above] {\small $x_2$};
\node at (3.5,-1.2) {\small Szenario 2};
\foreach \p in {(1.2,1.0),(1.5,1.3),(1.0,1.5),(1.8,1.1),(1.4,0.8),(1.7,1.6),(1.1,1.2),(1.6,0.9)}
  \fill \p circle (2pt);
\foreach \p in {(5.2,3.6),(5.5,3.9),(5.0,4.1),(5.8,3.7),(5.4,3.4),(5.7,4.2),(5.1,3.8),(5.6,3.5)}
  \fill \p circle (2pt);
\end{scope}
\end{tikzpicture}
\end{center}

Beurteilen Sie für beide Szenarien jeweils, ob (i) Elliptic Envelope / Mahalanobis,
(ii) KDE und (iii) GMM eine dichte Beschreibung der Normaldaten erwarten lassen.
Begründen Sie jede Entscheidung.

Szenario 1:
\antwortlinien{4}

Szenario 2:
\antwortlinien{4}

\section*{Aufgabe 6 -- EM-Algorithmus sortieren}

Bringen Sie die folgenden Schritte des GMM-Trainings in die richtige Reihenfolge
(1--5) und beschreiben Sie kurz, was im Expectation- und im Maximization-Schritt passiert:

\begin{itemize}
  \item[$\square$] Aktualisiere die Verteilungsparameter per Maximum-Likelihood anhand der
  Punkte mit den jeweils höchsten Wahrscheinlichkeiten (Maximization)
  \item[$\square$] Wähle Anzahl an Clustern bzw. Modi
  \item[$\square$] Berechne für jeden Datenpunkt die Wahrscheinlichkeit für jede Verteilung
  (Expectation)
  \item[$\square$] Wähle initiale Parameter pro Cluster (z.\,B. durch Vorlauf von k-Means)
  \item[$\square$] Wiederhole bis Konvergenz oder maximale Iterationszahl erreicht
\end{itemize}
\antwortlinien{3}

\vspace{10pt}
\begin{center}
\footnotesize\itshape
Nach Bearbeitung: Antworten an Claude schicken (Foto genügt) -- Korrektur mit Begründungs-Check.
\end{center}

\end{document}
